\subsection{Health and safety impacts}

Health is also a critical dimension in \acp{bss} promotion, complementing the environmental and transport outcomes discussed earlier~\cite{DeMaio2009, Fishman2016, Qiu2018}. This section explores the health and safety implications of bike-sharing, beginning with evidence of user-level physical and mental benefits, then reviewing the results of \acp{hia} studies that quantify population-level effects through physical activity, air pollution, and traffic-safety pathways, and concluding with a discussion of equity considerations and policy implications.

    \subsubsection{User-level health benefits of cycling}

    Regular cycling has been consistently associated with substantial health benefits, including improved cardiovascular fitness and lower risks of diabetes, cancer, and premature mortality~\cite{Green2021, Celis-Morales2017, Dutheil2020}. It provides an efficient means of achieving the recommended 150 minutes of moderate to vigorous physical activity per week~\cite{Knowler2002}, while also supporting mental health by reducing stress, alleviating symptoms of depression, and enhancing overall well-being~\cite{Chekroud2018, Ruiz2015}. By extending access to bicycles and encouraging regular use, \acp{bss} have the potential to translate these individual-level health benefits into broader population gains, but to what extent have such impacts been empirically demonstrated?

    Over the past decade, a growing body of research has examined how bike-sharing affects public health, particularly through changes in physical activity. Early observational and survey-based studies provided the first evidence of user-level health improvements. In 2010, \textit{BIXI Montréal} data showed that residents living near docking stations increased their cycling activity, suggesting that \acp{bss} can promote active travel~\cite{Fuller2013}. However, the study did not determine whether this increase came at the expense of other forms of active mobility. Around the same period, an eight-month study among university commuters in Valencia found that regular use of \textit{Valenbisi} contributed roughly half of the recommended 150 minutes of weekly moderate physical activity and was associated with modest reductions in BMI, highlighting a potential role for \acp{bss} in supporting healthy weight management~\cite{Molina-Garcia2013}. A survey of 3,000 \textit{Capital Bikeshare} users in Washington, D.C., similarly found that 31.5\% reported reduced stress and 30\% reported weight loss as a result of using the system~\cite{Alberts2012}. Complementing these self-reported outcomes, Xu~\cite{Xu2019} linked U.S. bike-sharing deployment data with county-level obesity prevalence statistics and found an average 1\% decline in obesity rates following \acp{bss} implementation, although this aggregate approach may have overestimated the effect due to spatial and behavioral heterogeneity~\cite{Hosford2019}.

    Expanding to a multi-city perspective, Fishman et al.~\cite{Fishman2014} analyzed systems in Melbourne, Brisbane, Washington D.C., London, and Minneapolis–St. Paul, reporting substantial overall increases in physical activity (from about 1.4 million additional minutes per year in Minneapolis–St. Paul to 74 million in London). However, the authors emphasized that the magnitude of health gains depends strongly on modal substitution: when bike-share trips replace car travel, users gain additional physical activity and related health improvements, whereas substituting walking trips produces the opposite effects. Similar results were reported for \textit{Dublinbikes}, where many trips replaced walking rather than motorized travel, leading to modest health improvements~\cite{Murphy2015}. Therefore, health benefits are contingent not only on \ac{bss} usage levels but also on the modes replaced.

    Taken together, these studies demonstrate that bike-sharing can effectively promote cycling and deliver measurable improvements in individual fitness, mental well-being, and weight management. Yet translating these user-level benefits into city-wide health gains remains challenging. A narrative review by Bauman et~al.~\cite{Bauman2017} found little evidence that current systems substantially increase overall physical activity for the general population, as most \ac{bss} users were already active and many trips substituted for walking or public transport rather than car travel. According to the authors, because regular users account for only a small share of the urban population, even if all met World Health Organization activity guidelines through cycling, the proportion of adults achieving these recommendations would rise by less than one percentage point. This suggests that significant health improvements at a city level are likely to occur only when bike-sharing reaches much broader segments of society.
        

    \subsubsection{\acp{hia} framework}

    While empirical studies provide valuable insights into individual behaviour and perceived health changes, they offer limited evidence of city-wide impacts or the specific pathways through which health effects occur. To address this gap, researchers have employed \acp{hia} as modelling frameworks to estimate the combined effects of bike-sharing on health outcomes. \acp{hia} estimate the net population-level health impacts by combining travel behavior data with epidemiological dose–response relationships~\cite{Rojas-Rueda2011, Woodcock2014, Otero2018, Clockston2021}. They typically evaluate changes in key outcomes such as mortality, morbidity, and Disability-Adjusted Life Years (DALYs), accounting for three principal pathways~\cite{Teixeira2021}: (i) health gains from increased physical activity, (ii) health risks from traffic injuries, and (iii) health losses from exposure to air pollution. This holistic approach allows researchers to study both the benefits and risks of \acp{bss} under different scenarios.

    To quantify physical activity, these models use the concept of the Metabolic Equivalent of Task (MET), a standardized physiological measure that expresses the energy expenditure of a specific activity as a multiple of the resting metabolic rate~\cite{Ainsworth2011}. One MET corresponds to the energy cost of sitting quietly, equivalent to an oxygen uptake of approximately 3.5~\si{\milli\liter}~\ce{O2} per kilogram of body weight per minute. Thus, an activity with an intensity of 6~METs requires roughly six times the energy expended at rest. Cycling intensity is typically assumed to range between 6 and 8~METs for conventional bicycles and around 6~METs for e-bikes with standard assistance~\cite{Rojas-Rueda2011, Woodcock2014, Otero2018, Clockston2021}. In this context, the total dose of physical activity is first estimated at the individual level, based on trip duration, frequency, and the proportion of the population using the \ac{bss}. This dose is expressed in MET-hours per week or per trip, reflecting the additional energy expenditure relative to rest. The resulting values are then combined with epidemiological dose–response functions derived from large cohort studies or meta-analyses to estimate the corresponding reductions in disease risk. Because these relationships are empirically non-linear, some \acp{hia} incorporate functions that assign larger relative benefits to previously sedentary populations and smaller marginal gains to already active individuals~\cite{Otero2018, Clockston2021, Chen2023}. Through this process, \acp{hia} translate individual-level physical activity increases associated with \acp{bss} into population-level estimates of prevented deaths or gains in DALYs.

    Air-pollution exposure is typically assessed in \acp{hia} through the inhalation of fine particulate matter (\ce{PM_{2.5}}) during travel, as this pollutant shows the strongest association with all-cause mortality and disease incidence~\cite{Woodcock2014, Otero2018, Clockston2021, Cao2021}. The inhaled dose is computed by combining background \ce{PM_{2.5}} concentrations with activity-specific ventilation rates, which are higher for cyclists than for pedestrians or motorists, and trip duration. This value represents the total pollutant mass inhaled per trip, which is then linked to health outcomes using the dose–response functions derived from epidemiological studies. Even though other pollutants such as \ce{NO2} or black carbon could also be modelled, most studies focus on \ce{PM_{2.5}} as they are highly correlated and produce similar health responses~\cite{Otero2018}. Moreover, for simplicity, assessments restrict the analysis to in-transit exposure and do not account for the indirect improvements in air quality resulting from reduced car use~\cite{Clockston2021}. 

    For context, the World Health Organization (WHO) recommends since 2021 long-term average \ce{PM_{2.5}} concentrations not exceeding 5~\si{\micro g/m^3} and daily averages below 15~\si{\micro g/m^3}~\cite{WHO2021}. However, air-quality levels in many cities, particularly in low- and middle-income countries, remain well above these thresholds. In Chinese megacities, for example, annual mean \ce{PM_{2.5}} concentrations on from 2004 to 2013 have ranged from 32 to 120~\si{\micro g/m^3}~\cite{Ma2016}, often exceeding international guidelines by an order of magnitude. Such elevated pollution levels imply that even modest increases in active travel can entail greater inhalation of pollutants, which may offset part of the health benefits from physical activity.

    Finally, traffic-safety impacts are measured by estimating the number of injuries and fatalities expected when substituting other travel modes with cycling. Models combine exposure metrics such as kilometres or hours cycled with mode-specific crash rates derived from police or hospital records~\cite{Otero2018, Clockston2021}. These rates are then applied to the estimated number of new cycling trips to calculate changes in injury and fatality risk relative to the baseline.

    \subsubsection{Empirical findings from HIAs}

    \acp{hia} have consistently found that the overall health impacts of \acp{bss} are driven primarily by gains in physical activity, while the air-pollution and traffic-injury pathways play comparatively minor roles. This section reviews empirical applications across different contexts, beginning with the main findings and the dominant influence of the physical-activity pathway, and then examining the air-pollution and traffic-injury pathways.


    \paragraph{Overview of \acp{hia} applications and the central role of the physical-activity pathway.}

    One of the first studies of this kind assessed Barcelona’s \textit{Bicing} program~\cite{Rojas-Rueda2011} in 2009. The analysis estimated that increased physical activity among users prevented approximately 12.5 deaths per year, while the combined risks from traffic crashes and pollution exposure were minimal by comparison. The net effect corresponded to 69 deaths avoided per million users annually, yielding a benefit–risk ratio of approximately 77:1. However, this outcome partly reflected an optimistic modeling assumption: the authors assumed that about 90\% of \textit{Bicing} trips replaced car journeys. This strong mode-shift assumption—acknowledged by the authors as a limitation—likely amplified the estimated health benefits relative to scenarios where trips would have replaced walking or public transport. A similar study conducted in 2010 in London’s \textit{Barclays Cycle Hire} incorporated empirical trip data and disaggregated outcomes by sex and age~\cite{Woodcock2014}. Using a comparative risk assessment approach, the authors found a smaller but still positive overall health benefit compared to Barcelona, with the magnitude of gains varying across demographic groups: men experienced larger benefits because they cycled more frequently and for longer distances, older adults benefited more in absolute terms due to higher baseline mortality risk reductions, while women faced higher relative injury risks in traffic.     
    
    Building on these city-level analyses, Otero et al.~\cite{Otero2018} assessed 12 large-scale \acp{bss} across Europe in 2018, estimating a combined 5.17 deaths avoided per 1,000 bicycles in service per year. The study explicitly included both \acp{m-b} and \acp{e-b}, distinguishing between standard- and high-assistance e-bikes. Although e-bikes provide lower levels of physical exertion per trip, the authors argue that they can still have a positive health impact, as they attract a wider range of users. Overall, the authors reported a benefit–risk ratio of 19:1, indicating that physical-activity gains outweighed risks from air pollution and traffic injuries by a wide margin. They also found that larger systems tend to generate greater aggregate health benefits, as the higher number of users and trips amplifies population-level effects. Although absolute numbers remained modest, the study highlighted that benefits could scale up substantially if a higher proportion of \acp{bss} trips replaced car journeys; for instance, increasing the share of car-substituted trips to 50\% could raise the number of deaths avoided nearly tenfold.
    
    In the United States, Clockston et al.~\cite{Clockston2021} conducted the first nationwide \ac{hia} of \acp{bss}, using 2019 ridership data to quantify both mortality and morbidity effects. Across all U.S. systems, BSS use was estimated to prevent approximately 4.7 premature deaths and 737 DALYs per year, equivalent to about \$36 million in avoided health costs. The authors also examined New York City separately, where BSS trips prevented around 2 premature deaths and 355 DALYs annually, yielding health savings of roughly \$15 million. Unlike the earlier European analysis by Otero et al.~\cite{Otero2018}, which focused solely on mortality, this assessment incorporated morbidity through disability-adjusted life years and economic valuation. The study further distinguished among substitution modes—car, walking, and public transport—showing that the largest health gains occurred when BSS replaced car trips, followed by walking, while replacing public-transport trips produced smaller or even neutral effects. Consistent with previous evidence, the increase in physical activity accounted for about 90\% of the total health impacts.

    Lastly, extending this evidence beyond Europe and North America, Chen et al.~\cite{Chen2023} analyzed Shanghai’s \textit{Mobike} program, offering the first \ac{hia} of a \ac{bss} in a highly polluted Asian megacity. Even under average \ce{PM_{2.5}} concentrations of about 45~\si{\micro g/m^3}, the study found that the gains from physical activity clearly dominated, preventing roughly 7–8 premature deaths per year among 306,000 users. Therefore, the overall health balance remained positive up to very high pollution levels. Only when \ce{PM_{2.5}} concentrations approached 68~\si{\micro g/m^3} did the pollution burden begin to offset the benefits. These results reaffirm the central role of physical activity as the main driver of health gains from \acp{bss}, even in environments with substantial air-quality challenges.


    \paragraph{Air-pollution pathway.}

    Across studies, the air-pollution pathway has been found to contribute only marginally to the overall health balance of \acp{bss}. In both Barcelona and London, the additional exposure to \ce{PM_{2.5}} during cycling was outweighed many times over by the benefits of increased physical activity, with the balance turning negative only under unrealistically high pollution scenarios~\cite{Rojas-Rueda2011, Woodcock2014}. Extending the analysis to 12 European cities, Otero et al.~\cite{Otero2018} similarly identified \ce{PM_{2.5}} as the least influential of the three modelled pathways but noted that none of the cities met the World Health Organization’s air-quality guidelines, suggesting that cleaner urban air would further amplify net benefits.
    
    
    Consistent with this pattern, Clockston et al.~\cite{Clockston2021} confirmed that pollution effects remain small at the national scale, estimating less than one death attributable to in-transit \ce{PM_{2.5}} exposure per year. Their analysis also showed that the health impact of air pollution depends strongly on which modes are replaced. A counterintuitive result emerged when bike-share trips substituted for walking: exposure to \ce{PM_{2.5}} slightly decreased because cycling trips were shorter in duration, leading to lower total inhaled volumes despite higher ventilation rates. By contrast, when trips replaced car or public transport, exposure increased, as cyclists breathe more intensely and are directly immersed in the traffic stream~\cite{Rojas-Rueda2011, Otero2018, Clockston2021}. This asymmetry highlights that replacing sedentary modes is not always beneficial from an air-quality perspective, even if it yields large gains in physical activity. Overall, however, the evidence consistently shows that these pollution-related risks remain minor compared with the substantial health benefits derived from increased physical activity.

    Nevertheless, a few studies have shown that the air-pollution pathway can play a more prominent role under specific conditions. Chen et al.~\cite{Chen2023} reported that in Shanghai’s \textit{Mobike} program, elevated \ce{PM_{2.5}} concentrations significantly affected the overall health balance. Under average levels of about 45~\si{\micro g/m^3}, the benefits from physical activity still outweighed the harms from pollution, but as concentrations approached 68~\si{\micro g/m^3}, the air-pollution burden began to offset these gains, eventually reversing the net outcome. This finding highlights that in highly polluted cities, the contribution of the air-pollution pathway can become decisive, potentially turning otherwise beneficial cycling activity into a net health loss.
    
    On the other hand, Hou et al.~\cite{Hou2022} found a positive air-quality balance in their assessment of New York City’s \textit{Citi Bike} program. The authors estimated an overall benefit of 5,745 DALYs and 216 premature deaths avoided, with major contributions stemming from reductions in \ce{PM_{2.5}} emissions (-2,757 DALYs), alongside additional benefits from lower traffic injuries and smaller gains from increased physical activity (-10 DALYs). The prominence of the air-pollution pathway reflects both methodological and contextual differences: their model accounted for system-wide emission reductions resulting from decreased car use, whereas most earlier \acp{hia} considered only users’ in-transit exposure. Moreover, New York’s high baseline walking rate of approximately 35\% limited the additional physical-activity gains achievable through cycling. A sensitivity analysis further showed that health improvements were greatest when bike-share trips replaced car travel, while substitution from walking produced smaller or even negative changes in physical-activity-related DALYs.

    Beyond \acp{hia} frameworks and their traditional reliance on daily mean \ce{PM_{2.5}} values, recent studies have also underscored the importance of short-term peaks and behavioural adaptation in shaping cyclists’ real-world exposure. Cao et al.~\cite{Cao2021} proposed two novel exceedance-based indicators—Exceedance Exposure Risk and Accumulative Exceedance Exposure Risk—that integrate minute-level pollution data with individual trip records to capture transient exposure surges. Their analysis of more than one million \textit{Mobike} trips in Guangzhou revealed that riders were exposed to markedly higher \ce{PM_{2.5}} risks during evening traffic peaks, when exceedance concentrations clustered around dense commercial and leisure areas. This approach demonstrated that daily averages can underestimate inhaled doses by up to an order of magnitude and underscored the need for high-resolution spatiotemporal modelling in future \acp{hia}. Complementing this methodological refinement, Park et al.~\cite{Park2022} analysed three years of data from Seoul’s public bike-sharing system and found a clear non-linear response between ambient \ce{PM_{2.5}} and ridership: usage declined sharply once concentrations exceeded roughly 38–42~\si{\micro g/m^3}. Commuting trips showed pronounced risk-averting behaviour, whereas leisure trips remained less sensitive to pollution levels.


    \paragraph{Traffic-safety pathway.}

    The traffic-injury pathway also contributes far less to total health outcomes than the physical-activity gains but shows greater variability across cities. In general, substituting car trips with cycling increases fatality risk in cities lacking protected lanes or safe junctions~\cite{Otero2018}. Yet this risk can be substantially mitigated where infrastructure is well designed. Among the European cases analyzed by Otero et al.~\cite{Otero2018}, \textit{Seville} was the only city where car-to-bike substitution reduced fatality risk, owing to its rapid expansion of segregated lanes and traffic-calming measures. A similar pattern emerged in the United States, where nationwide modeling found that, when all substitutions were considered together, the overall shift to \acp{bss} slightly increased users’ risk of traffic incidents, yet replacing walking trips actually reduced it~\cite{Clockston2021}. This counterintuitive outcome arises because walking exposes users to potential pedestrian–vehicle conflicts for longer periods, while cycling covers the same distance more quickly. Several mitigating mechanisms have also been identified: London’s experience and other observational evidence suggest that the design of shared bicycles—being heavier, slower, and more visible—may contribute to lower crash severity and reduced injury rates~\cite{Woodcock2014}. Furthermore, the well-documented “safety-in-numbers” effect implies that as cycling volumes grow, per-cyclist collision risk declines due to heightened driver awareness~\cite{Jacobsen2015}. Overall, although traffic-related risks can offset a portion of the health gains in specific contexts, they remain modest compared with the benefits of increased physical activity—particularly in cities that invest in safe, well-designed cycling infrastructure.

    Safety outcomes represent another critical dimension of user-level impacts. A systematic review by Chen et~al.~\cite{Chen2020} synthesized evidence from multiple cities and found that bike-share users are substantially less likely to wear helmets than private cyclists, with unhelmeted rates typically ranging between 36\% and 89\%. While changes in injury incidence following the introduction of \acp{bss} were mixed across contexts, the review highlighted that low helmet use remains a consistent pattern among casual and first-time riders. Martin et~al.~\cite{Martin2016} similarly concluded that safety effects vary widely between cities: some reported small increases in crashes, others decreases, often depending on local infrastructure, traffic conditions, and cycling culture. However, across most settings, per-trip injury rates remained low, and several studies supported the “safety-in-numbers” hypothesis whereby growing cyclist volumes reduce per-capita crash risk. Together, these findings suggest that although \acp{bss} can encourage new and less experienced riders, their overall safety performance is comparable to, or in some cases better than, that of private cycling when supported by adequate infrastructure and visibility measures.

    Safety outcomes represent the other critical dimension within HIAs. Helmet use is one of the most frequently discussed issues: bike-share riders are less likely to wear helmets than regular cyclists, often because casual or impromptu trips make carrying a helmet inconvenient~\cite{Fishman2016}. Studies in North America documented very low helmet-wearing rates among bike-share users (often under 15\%). Although the impact of this on overall head injuries is debated, a study by Graves et al.~\cite{Graves2014} found that the proportion of head injuries among all bicycle-related injuries increased in cities after implementing bike-share programs without mandatory helmet laws. This suggests that while the absolute number of injuries did not spike, a higher fraction of injured cyclists were helmetless bike-share users—a pattern warranting attention. Some cities have considered or implemented measures such as helmet rental kiosks or education campaigns to improve helmet uptake, though these sometimes conflict with the spontaneity that makes bike-share attractive. Another safety aspect is the exposure of new or inexperienced cyclists to busy urban traffic. Bike-share often entices people who might not otherwise cycle to give it a try; ensuring that those users have safe cycling infrastructure (protected lanes, traffic-calmed routes) is crucial to minimizing crashes. Thankfully, many bike-share programs coincide with city efforts to improve cycling conditions more broadly. 

    % Complementary empirical evidence from post-implementation analyses shows that shared bicycles often exhibit lower crash rates per trip than private cycling~\cite{Martin2016}, though the proportion of head injuries may rise due to low helmet use~\cite{Chen2020}.
    
    % Moreover, their results lead to the conclusion that if anything bike sharing may be safer than regular cycling~\cite{Woodcock2014}.

    %Otero2018: traffic fatalities can be influenced by multiple causes, traveler behavior, infrastructure, traffic laws, and mode of transport. In all the cities (with the exception of Seville) we found that car-bike substitution increases the risk of traffic fatalities in the travelers (see Fig. 2). Milan, Brussels, and Warsaw provided the highest risk of traffic fatality between cities when the car is substituted by bike. Seville is the only city that reported a lower risk of traffic fatalities in bike compared to a car. This can be explained because in the last few years Seville has invested in bike infrastructure, especially in segregated bike lanes, traffic signaling, and bike promotion and education

    % Clockston2021: In this study, traffic incidents were modeled to estimate the risk of traffic injuries and fatalities among BSS users. In both NYC and the U.S,the current substitution from all modes of transport to BSS overall was estimated to increase the risk of traffic incidents (injuries and fatalities)among BSS users. However, when broken down according to mode, substituting walking for BSS trips decreases the risk of traffic incidents in both NYC and the U.S. (see tables S27 and S57 in the supplemental material). In this study, we did not analyze the impact of specific infrastructure, road and weather conditions, and other local factors that are known to have an impact on traffic safety (I. Otero et al., 2018).

    % Teixeira2021
    % Conversely, the exposure of BSS users to traffic crashes and air pollution could negatively offset health gains. Nevertheless, the current evidence reveals these risks to beon the scale of cycling in general (Fishman & Schepers, 2016; Fuller, Gauvin, Morency,Kestens, & Drouin, 2013; Woodcock et al., 2014). For instance, Fuller, Gauvin, Morency,Kestens, et al. (2013) investigated the traffic risks of BIXI users (Montreal’s BSS) throughthree cross-sectional surveys (at launch and at the end of the first and second seasonsof BIXI). The study indicated that there was no change in the likelihood of reported col-lisions or near misses, even though the city’s cycling levels had increased (Fuller,Gauvin, Morency, Kestens, et al., 2013). These results seem to be in line with the safetyin numbers hypothesis, where an increase in the number of cyclists does not lead to aproportional increase in the number of injuries (Fuller, Gauvin, Morency, Kestens, et al.,2013; Jacobsen, 2003). Some data seems to point that the injury risk may be evenlower for BSS users (Fishman & Schepers, 2016; Woodcock et al., 2014). Fishman and Sche-pers (2016), comparing hospital injury data of ten North American cities (five with BSS andfive as control), revealed a decrease on cyclists’injury risk after the BSS implementation,with the risk decrease large enough to outbalance the increase on cycling use. Possible explanations include lower speeds or more awareness among car drivers (Fishman &Schepers, 2016), but this trend requires further research.

    % Importantly, the safety record of station-based systems in the U.S. has been strong. From 2007 to 2014, there were no fatal crashes involving bike-sharing bicycles, and crash rates have generally been lower than for private bicycle use. These findings are consistent with the “safety in numbers” hypothesis, which suggests that as cycling volumes rise, driver awareness improves and per-cyclist collision risks fall~\cite{Jacobsen2015}.

    \subsubsection{Equity and policy implications}

Finally, several studies have begun examining how health outcomes vary across population groups and urban contexts. In New York City, Babagoli et al.~\cite{Babagoli2019} showed that the distribution of BSS infrastructure in New York City disproportionately favored wealthier neighborhoods, which in turn affected who received the health benefits. Similar analyses emphasize the need for disaggregated data to assess equity implications by income, gender, and age~\cite{Clockston2021}. Despite methodological advancements, HIAs of \acp{bss} still face several challenges. Results are sensitive to assumptions about the modes replaced by BSS use: studies that assume a high share of car replacement tend to report larger health gains than those where walking or transit are substituted. There is also ongoing debate about how to model long-term behavior change (e.g., whether using a \ac{bss} initiates lifelong cycling habits), which could significantly increase the estimated health benefits~\cite{Bauman2017}. Furthermore, most assessments focus on physical health, with limited treatment of mental health outcomes or chronic morbidity beyond cardiovascular disease. \textcolor{red}{It is also important to note that the evidence base remains concentrated in high-income countries with lower pollution levels and stronger traffic safety standards~\cite{Teixeira2020, WHO2018}. In cities with higher air pollution or weaker safety conditions, the balance between benefits and risks may differ substantially~\cite{Tainio2016}.} Finally, data limitations—especially in lower-income or smaller cities—restrict the scalability and generalizability of many findings.


In summary, the collective evidence shows that bike-sharing produces clear net health gains, driven primarily by increased physical activity. Comprehensive assessments indicate that these gains greatly outweigh the relatively small health losses from air pollution and crash risks. Nevertheless, future research should examine long-term behavioral change, mental health outcomes, and equity dimensions. Overall, \acp{bss} represent a cost-effective public health intervention~\cite{Yu2018}, particularly when designed to replace car travel, protect vulnerable users, and ensure equitable access to cycling opportunities.



%Fishman2016 té un apartat 5.2. health i 5.3. Road traffic injury and bss

% ChatGPT
% *These results are consistent with the “safety in numbers” hypothesis, which posits that as cycling volumes grow, driver awareness improves and the probability of collision per cyclist decreases~\cite{Fishman2016, Fuller2013, Jacobsen2003}.* This may also reflect the upright design, moderate speed, and central location of most bike-share operations~\cite{Shaheen2014}..